\chapter{Versione 2}

\section{Obbiettivi e Requisiti}
La seconda versione dell'applicazione introduce la possibilità per il Configuratore di proporre vere e proprie iniziative basate sulle Categorie precedentemente create e di renderle pubbliche sulla nuova \textbf{Bacheca}.

\section{Casi d'Uso (Aggiunte rispetto a V1)}

I casi d'uso base del Configuratore restano invariati. A essi si aggiungono le seguenti funzionalità:

\begin{enumerate}
    \setcounter{enumi}{7}
    \item \textbf{Creare Proposta}: Il configuratore seleziona una categoria e compila un modulo per proporre un nuovo evento. Il sistema verifica la congruenza logica delle date e la compilazione dei campi obbligatori, assegnando lo stato \texttt{VALIDA}.
    \item \textbf{Pubblicare Proposta in Bacheca}: Il configuratore sceglie di rendere pubblica una proposta valida, facendola passare allo stato \texttt{APERTA} nella bacheca di sistema.
    \item \textbf{Visualizzare Bacheca}: Il configuratore può vedere tutte le proposte attualmente aperte, raggruppate per categoria d'appartenenza.
\end{enumerate}

\subsection{Diagramma UML dei Casi d'Uso (V2)}
% Inserire qui l'immagine generata dal PlantUML
Le aggiunte specifiche di questa versione (Creazione Proposta, Pubblicazione, e Visualizzazione Bacheca) sono state evidenziate rispetto al diagramma V1.

\begin{figure}[h]
    \centering
    % \includegraphics[width=0.8\textwidth]{images/uc_v2.png}
    \caption{Diagramma dei Casi d'Uso - Versione 2}
    \label{fig:uc_v2}
\end{figure}

\section{Diagramma delle Classi (Aggiunte rispetto a V1)}

Sono state introdotte le nuove entità \texttt{Proposta}, \texttt{StatoProposta} e la \texttt{Bacheca}, oltre al controller dedicato \texttt{GestoreProposte}.
% Inserire qui l'immagine generata dal PlantUML
\begin{figure}[h]
    \centering
    % \includegraphics[width=1\textwidth]{images/class_v2.png}
    \caption{Diagramma delle Classi (Delta evidenziati) - Versione 2}
    \label{fig:class_v2}
\end{figure}
